\section{从图形学到机器人：仿真里的控制如何进入真实世界}

本章从谭杰的研究路径开始。他博士做计算机图形学，尤其是 physics-based character animation，即在仿真环境里让人形或动物形态自然行走。这个方向和机器人高度相似：图形学是在 simulation 里做机器人，机器人是在真实世界里做图形学。区别在于，仿真里可以拿到所有状态信息，真实世界则充满噪声、摩擦和不确定性。

\begin{figure}[H]
\centering
\includegraphics[width=0.96\textwidth]{graphics-to-robotics.png}
\caption{从图形学到机器人：在仿真里做机器人，再迁移到真实世界。自制概念图，依据 00:02:00--00:13:06 对谈内容整理。}
\end{figure}

\begin{knowledgebox}{读图：图形学提供了“先在仿真里学会”的路线}
图中从图形学和物理动画开始，经由强化学习学习步态，再通过 Sim2Real 迁移到真实机器人。谭杰早期工作的关键，就是把仿真里的控制能力带到四足机器人真实运动中。
\end{knowledgebox}

\subsection{Sim2Real 与强化学习的第一次范式转移}

本节先看“身体怎么动”的问题。大模型出现之前，机器人领域已经发生过一次重要转向：让机器人先在仿真里通过强化学习学会运动，再迁移到真实硬件。

Sim2Real 是 simulation-to-real 的缩写，即先在仿真环境中训练策略，再迁移到真实机器人。谭杰提到他在 Google 的早期工作 Sim2Real Learning Agile Locomotion for Quadruped Robots，用深度强化学习解决四足机器人敏捷运动。强化学习是让智能体在环境中试错，通过 reward 信号改进行为的训练方法；PPO 是其中一种常用算法。

\begin{importantbox}{第一次范式转移}
过去机器人运动控制高度依赖传统控制和 Model Predictive Control；过去五到十年，强化学习 + Sim2Real 基本改变了步态和 locomotion 领域，让跑、跳、翻、打拳等运动能力快速普及。
\end{importantbox}

\subsection{大模型带来的第二次范式转移}

接下来再看“机器人是否懂任务”的问题。运动控制解决的是身体，大模型补上的则是语言、常识和计划，这也是为什么谭杰把大模型比作大脑。

大模型出现之前，很多机器人没有 common sense，也不懂自然语言。你要机器人做咖啡，必须写程序或拆成非常具体的控制指令。大模型带来的改变是：机器人可以理解自然语言，拆解任务步骤，知道“做咖啡”大概需要杯子、水、咖啡、加热和倒入等常识。谭杰用大脑和小脑类比：大模型更像大脑，负责思维、语言、计划；强化学习更像小脑，负责平衡、运动和执行。

\begin{figure}[H]
\centering
\includegraphics[width=0.94\textwidth]{brain-cerebellum-robot.png}
\caption{大脑与小脑：多模态模型负责认知，强化学习负责执行。自制概念图，依据 00:02:00--00:13:06 对谈内容整理。}
\end{figure}

\begin{knowledgebox}{读图：机器人需要大脑，也需要小脑}
大脑负责语言理解、常识和任务规划；小脑负责步态、平衡、抓取和底层控制。只有大脑，机器人会“知道该做什么但做不好”；只有小脑，机器人会“能动但不懂目标”。
\end{knowledgebox}

\subsection{本章小结}

谭杰的研究路径解释了机器人过去十年的两个转折：第一，强化学习和 Sim2Real 改变了运动控制；第二，大模型把语言、常识和计划带入机器人。今天的机器人研究要同时处理大脑和小脑。

